\documentclass[10pt,a4paper]{amsart}
\usepackage[margin=19mm]{geometry}
\usepackage{amsmath,amssymb,mathtools}
\usepackage[scr=rsfs,cal=euler]{mathalfa}
\usepackage{xfrac}
\usepackage[unicode,hidelinks,bookmarks=false]{hyperref}
\newcommand{\ie}{i.e.\ }
\newcommand{\cf}{cf.\ }
\newcommand{\resp}{respectively\ }
\newcommand{\Subsection}{\subsection}
\providecommand{\nobreakdash}{\nobreak\hbox{-}}
\setlength{\emergencystretch}{2em}
\multlinegap0pt
\usepackage{bm}
\usepackage{bbm}
\usepackage{dsfont}
\usepackage{xspace}
\usepackage{cancel}
\usepackage{mathtools}
\usepackage{amsthm}
\makeatletter
\@ifundefined{lem}{\newtheorem{lem}{Lemma}}{}
\@ifundefined{thm}{\newtheorem{thm}{Theorem}}{}
\makeatother
\newcommand{\C}{\mathbb{C}}
\newcommand{\Hfr}{\mathfrak{H}}
\newcommand{\N}{\mathbb{N}}
\newcommand{\R}{\mathbb{R}}
\newcommand{\Ss}{\mathbb{S}}
\newcommand{\Ub}{\mathbf{U}}
\newcommand{\be}{\begin{equation}}
\newcommand{\ee}{\end{equation}}
\newcommand{\ii}{\mathrm{i}}
\newcommand{\ov}{\overline}
\newcommand{\rank}{\mathrm{rank}}
\newcommand{\sv}{\mathbf{s}}
\newcommand{\uu}{\mathbf{u}}
\renewcommand{\geq}{\geqslant}
\renewcommand{\leq}{\leqslant}
\title[Global Well-Posedness and Soliton Resolution for the Half-Wa]{Global Well-Posedness and Soliton Resolution for the Half-Wave Maps Equation with Rational Data}
\author{Patrick G\'erard, Enno Lenzmann}
\date{Source: arXiv:2412.03351v3 (CC BY 4.0)}
\begin{document}
\maketitle

\noindent\emph{Source and licence.} Global Well-Posedness and Soliton Resolution for the Half-Wave Maps Equation with Rational Data. Version arXiv:2412.03351v3 (CC BY 4.0), distributed under the Creative Commons Attribution 4.0 International licence, as recorded in the arXiv licence metadata for this exact version and shown on the abstract page https://arxiv.org/abs/2412.03351v3. Full rights notice: licenses/arxiv-2412.03351-CC-BY-4.0.txt.\par\smallskip

\noindent\textit{Editorial scope.} Counted unit: the Thm whose source label is \texttt{thm:toeplitz\_stereo\_app} in the article (source0.tex lines 2582--2596, source label \texttt{thm:toeplitz\_stereo\_app}), delivered together with the article's own complete original proof of it (source0.tex lines 2600--2737, one \texttt{proof} environment of 138 lines). No mathematics is written, repaired or re-derived: statement and proof are the article's own text line for line. The proof is detached from the statement in the source: the article states the result at lines 2582--2596 and prints its proof at lines 2600--2737, and the proof environment's own header names the statement, so the association is the article's own. Required context retained uncounted: lem:R\_n (lines 2096--2098); lem:simple\_poles (lines 2497--2503); lem:mult\_poles (lines 2546--2555). Every definition, display and label of those lines is retained unchanged. Omitted from this package and not redistributed: 1--2095, 2099--2496, 2504--2545, 2556--2581, 2597--2599, 2738--3386. Nothing inside a retained excerpt is omitted or altered: every retained body is delivered exactly as the article's lines are written, so no omission receipt of any kind accompanies this package. Citations: the retained excerpts carry no citation command at all, so no bibliography entry of the article is transcribed and no reference is redistributed. Reference adaptation: none was needed and none was made; every reference inside the delivered unit resolves inside it, so no unresolved reference or citation is produced. Printed slips kept literal: no printed word, symbol, hypothesis or number of the counted statement or of its proof was altered. Layout and class: the article's own class is \texttt{amsart} with packages including amsmath, amssymb, amsthm, mathtools, dsfont, bm, bbm; the wrapper is the lane's pinned \texttt{amsart} editorial wrapper at 10pt on A4, which restates the theorem-like environments the retained text needs and adds the article's own macro definitions that the retained text needs (\texttt{\textbackslash C}, \texttt{\textbackslash Hfr}, \texttt{\textbackslash N}, \texttt{\textbackslash R}, \texttt{\textbackslash Ss}, \texttt{\textbackslash Ub}, \texttt{\textbackslash be}, \texttt{\textbackslash ee}, \texttt{\textbackslash geq}, \texttt{\textbackslash ii}, \texttt{\textbackslash leq}, \texttt{\textbackslash ov}, \texttt{\textbackslash rank}, \texttt{\textbackslash sv}, \texttt{\textbackslash uu}), with extra wrapper packages (bm, bbm, dsfont, xspace, cancel, mathtools). The delivered numbering is the unit's own. Assets: no retained span contains a figure environment, a graphics-inclusion command, a PDF-page inclusion, a table or a TikZ picture; no image file is copied into this package, the package declares no asset dependency and no image file is redistributed with it. Licence: version arXiv:2412.03351v3 of this article is distributed under the Creative Commons Attribution 4.0 International licence, as recorded in the arXiv licence metadata for this exact version and shown on the abstract page https://arxiv.org/abs/2412.03351v3. A full rights notice is delivered at licenses/arxiv-2412.03351-CC-BY-4.0.txt. Characters: every retained excerpt is pure ASCII, so the delivered engine, XeLaTeX with its default Latin Modern text font, reports no missing character.
\begin{lem} \label{lem:R_n}
The set $\mathcal{R}_N \subset \C(X)$ can be canonically identified with a non-empty, open and connected subset $\mathcal{A}_N$ in $\C^{2N}$. 
\end{lem}

\begin{lem} \label{lem:simple_poles}
Assume $\uu : \R \to \Ss^2$ is a rational function of the form
$$
\uu(x) = \uu_\infty + \sum_{j=1}^N \left ( \frac{\sv_j}{x-z_j} + \frac{\ov{\sv}_j}{x-\ov{z}_j} \right )
$$
with some integer $N \geq 1$, $\uu_\infty \in \Ss^2$, $\sv_1, \ldots, \sv_N \in \C^3 \setminus \{ 0 \}$, and pairwise distinct poles $z_1, \ldots, z_N \in \C_-$. Then it holds $\mathrm{Rank} (K_{\Ub}) = N$.
\end{lem}

\begin{lem} \label{lem:mult_poles}
Suppose that $\uu : \R \to \Ss^2$ is of the form
$$
\uu(x) = \uu_\infty + \sum_{j=1}^p \sum_{k=1}^{m_j} \left ( \frac{\sv_{j,k}}{(x-z_j)^k} + \frac{\ov{\sv}_{j,k}}{(x-\ov{z}_j)^k} \right )
$$
with some integers $N \geq 1$, $1 \leq p,m_j \leq N$, vectors $\sv_{j,k} \in \C^3 \setminus \{ 0 \}$, and pairwise distinct $z_1, \ldots, z_N \in \C_-$. Then it holds that
$$
\mathrm{Rank}(K_\Ub) \geq N = \sum_{j=1}^p m_j \, .
$$
\end{lem}

\begin{thm} \label{thm:toeplitz_stereo_app}
Let $\uu =(u_1, u_2, u_3) : \R \to \Ss^2$ be a rational map. Given an integer $N \geq 1$, the following statements are equivalent.
\begin{itemize}
\item[(i)] $\dim \Hfr_1 = \rank(K_{\Ub}) = N$.
\item[(ii)] The least common denominator of $u_1, u_2, u_3$ has degree $2N$.
\item[(iii)] There exists a rational function $R \in \C(X)$ of the form
$$
R(x) = \frac{P(x)}{Q(x)} \, ,
$$
where $P \in \C[X]$ is a polynomial of degree $N$ and $Q \in \C[X]$ is a non-zero polynomial of degree at most $N-1$, such that $P$ and $Q$ have no common factor such that, up to rotation on the sphere $\Ss^2$, we have
$$
u_1(x) + \ii u_2(x) = \frac{2 R(x)}{R(x) \ov{R}(x)+1} \, , \quad u_3(x) = \frac{R(x) \ov{R}(x) -1}{R(x) \ov{R}(x) + 1} \, .
$$
\end{itemize}
\end{thm}  

\begin{proof}[Proof of Theorem \ref{thm:toeplitz_stereo_app}]
We divide the proof into the following steps.

\medskip
\textbf{$(ii) \Rightarrow (iii)$}. Assume $\uu = (u_1, u_2, u_3) : \R \to \Ss^2$ is a rational map with the least common denominator given by a polynomial $D \in \R[X]$ of degree $2N$. (Note that $D$ must have even degree, since $u_1, u_2, u_3$ are real-valued rational functions with no poles in $\R$.) Moreover, up to a rotation on $\Ss^2$, we may assume that $u_3(x) \to 1$ as $|x| \to \infty$, so that there exist polynomials $Q_j \in \R[X]$ such that
$$
u_j(x) = \frac{Q_j(x)}{D(x)} \quad \mbox{for $j=1,2,3$},
$$
where $Q_1, Q_2$ have degree at most $2N-1$ and $Q_3$ has degree $2N$, with the same leading coefficient as $D$. Now the condition $u_1^2 + u_2^2 + u_3^2 = 1$ means that
$$
Q_1^2 + Q_2^2 + Q_3^2 = D^2 \, , 
$$
or equivalently
$$
(Q_1 + \ii Q_2)(Q_1 - \ii Q_2) = (D-Q_3)(D+Q_3) \, .
$$
Since $Q_3$ and $D$ have the same leading coefficient, the degree of $D+Q_3$ is $2N$ and the degree $\delta$ of $D-Q_3$ is at most $2N-1$. Denote by $d$ the degree of $Q_1 + \ii Q_2$. Since $Q_1$ and $Q_2$ are real polynomials, $d$ is also the degree of $Q_1 - \ii Q_2$ and hence
$$
2d = \delta + 2N \, .
$$
This implies
$$
N \leq d \leq 2N-1 \, .
$$
Furthermore, we recall that $D$ is the least common denominator of $u_1, u_2, u_3$ which means that $Q_1, Q_2, Q_3, D$ have no common factor, or equivalently the polynomials
$$
Q_1+ \ii Q_2, Q_1 - \ii Q_2, D-Q_3, D + Q_3
$$
have no common factor.

Now, we claim that $Q_1 + \ii Q_2$ and $D-Q_3$ have at least $d-N$ common zeros -- counted with multiplicities. Indeed, assume that $\alpha \in \C$ is a zero of $D-Q_3$ of multiplicity $m \geq 1$. We distinguish the following cases depending whether $\alpha \in \R$ or $\alpha \not \in \R$.

If $\alpha$ is real, then $\alpha$ is in fact a zero of $Q_1$ and $Q_2$, hence it is a zero of $Q_1+\ii Q_2$ and $Q_1 -\ii Q_2$ with the same multiplicity $\mu$. Since $\alpha$ cannot be a zero of $D+Q_3$ (otherwise $Q_1+\ii Q_2, Q_1 - \ii Q_2, D-Q_3, D+ Q_3$ would have common factor), we infer that $2 \mu = m$. 

If $\alpha$ is not real, then $\alpha$ is a zero of $Q_1 + \ii Q_2$ or $Q_1 - \ii Q_2$. Since $\ov{\alpha}$ is also a zero of the real polynomial $D-Q_3$, we can choose the zero $\beta \in \{ \alpha, \ov{\alpha} \}$ having the maximal multiplicity $\mu$ as zero of $Q_1 + \ii Q_2$. This shows $\mu \geq \frac{m}{2}$.

Summing up, we have found a common factor of $D-Q_3$ and $Q_1 + \ii Q_2$ with degree at least equal to half of the degree of $D-Q_3$, namely $d-N$. Therefore we can write
\be \label{eq:PQ}
\frac{Q_1 + \ii Q_2}{D-Q_3} = \frac{P}{Q}
\ee
where $P$ and $Q$ are polynomials in $\C[X]$ with no common factor, and $P$ has degree $d-r$, $Q$ has degree $2(d-N)-r$ for some $r \geq d-N$. Notice that $\deg P > \deg Q$.

Next, we prove that equality $r=d-N$ holds. Indeed, from \eqref{eq:PQ}, we conclude
$$
\frac{Q_1}{D} = \frac{P \ov{Q} + \ov{P} Q}{P \ov{P}+ Q \ov{Q}}, \quad \frac{Q_2}{D} = \frac{P \ov{Q}- \ov{P}Q}{\ii(P \ov{P} + Q \ov{Q})}, \quad \frac{Q_3}{D} = \frac{P \ov{P} - Q \ov{Q}}{P \ov{P}+ Q \ov{Q}} \, .
$$
This implies that $u_1, u_2, u_3$ have a common denominator of degree equal to $2 \deg P$. Hence
$$
2(d-r) \geq 2N \, ,
$$
which implies $r \leq d-N$, leading to the desired equality $r=d-N$. By defining the rational function
$$
R(x) = \frac{P(x)}{Q(x)}  \in \C(X),
$$
we conclude that (iii) holds. This completes the proof of the implication $(ii) \Rightarrow (iii)$.

\medskip
\textbf{$(iii) \Rightarrow (ii)$.} Suppose we are given a rational function
$$
R(x) = \frac{P(x)}{Q(x)} \, ,
$$
where $P$ is a polynomial of degree $N$, $Q$ is a polynomial of degree at most $N-1$, and $P, Q$ have no common factor. The formulae
$$
u_1 = \frac{R + \ov{R}}{R \ov{R}+1}, \quad u_2 = \frac{R - \ov{R}}{\ii (R \ov{R}+1)}, \quad u_3 = \frac{R \ov{R}-1}{R \ov{R}+1} 
$$
clearly define a rational map $\uu=(u_1,u_2, u_3)$ from $\R$ with values in $\Ss^2$. Furthermore, we see that $|P|^2 + |Q|^2$ is a common denominator of $u_1,u_2,u_3$ and is degree is $2N$. Let us prove that $|P|^2 + |Q|^2$ is the least common denominator of $u_1, u_2, u_3$. We argue by contradiction. Suppose there is a common factor of the polynomials
$$
P \ov{Q} + \ov{P} Q, P \ov{Q}- \ov{P} Q, P \ov{P} - Q \ov{Q}, P \ov{P}  + Q \ov{Q}
$$
or equivalently of the polynomials
$$
P \ov{Q}, \ov{P} Q, P \ov{P}, Q \ov {Q} \, .
$$
Since $P, Q$ have no common factor, there exist polynomials $U, V$ such that
$$
UP + VQ = 1 \, .
$$
Therefore the polynomials
$$
\ov{U} (P \ov{P} )+ \ov{V}(P \ov{Q}) = P, \quad \ov{U} (\ov{P} Q) + \ov{V} ( Q \ov{Q} ) = Q
$$
would have a common factor, which yields a contradiction. This proves that (iii) implies (ii).

\medskip
\textbf{$(ii) \Rightarrow (i).$} Let us assume that $\uu = (u_1, u_2, u_3) : \R \to \Ss^2$ has a least common denominator of degree $2N$. We claim that
\be \label{eq:rank_K}
\mathrm{Rank} (K_\Ub) = N \, .
\ee
Indeed, if $\uu : \R \to \Ss^2$ has only simple poles (i.\,e., the assumptions of Lemma \ref{lem:simple_poles} are satisfied), we can directly apply Lemma \ref{lem:simple_poles} to conclude that \eqref{eq:rank_K} holds.

To deal with the case multiple poles occurring in $\uu : \R \to \Ss^2$, we need the following approximation result.

\begin{lem} \label{lem:rat_approx}
Let $P \in \C[X]$ be a polynomial of degree $N \geq 1$, $Q \in  \C[X]$ be a non-zero polynomial of degree at most $N-1$, and assume that $P,Q$ have no common factor. Then there exist sequence $P_n, Q_n \in \C[X]$ such that $P_n, Q_n$ have no common factor and
$$
\deg P_n = N, \quad \deg Q_n \leq N-1, \quad P_n \to P, \quad Q_n \to Q \quad \mbox{in $\C[X]$} \, .
$$
Furthermore, the zeros of $|P_n|^2 + |Q_n|^2$ are simple for every $n \in \N$.
\end{lem}

\begin{proof}[Proof of Lemma \ref{lem:rat_approx}]
Consider the set $\mathcal{A}$ of pairs of polynomials $(P,Q) \in \C[X] \times \C[X]$ such that $\deg = N$, $\deg Q \leq N-1$ and $P,Q$ have no common factor and $P$ is monic. By Lemma \ref{lem:R_n}, we can identify $\mathcal{A}$ with a connected open subset in $\C^{2N}$. On the set $\mathcal{A}$, the condition that the discriminant of $|P|^2 + |Q|^2$ is different from 0 is an open dense subset. This completes the proof. 
\end{proof}

Suppose now that $\uu : \R \to \Ss^2$ has multiple poles and the least common denominator of $u_1, u_2, u_3$ has degree $2N$. By Lemma \ref{lem:mult_poles}, we must have
$$
\mathrm{Rank} \, (K_{\Ub}) \geq N \, .
$$

On the other hand, by the proven implication $(ii) \Rightarrow (iii)$, there exists a rational function 
$$
R(x) = \frac{P(x)}{Q(x)} \in \C(X)
$$
with $\deg P = N$, $\deg Q \leq N-1$ with $Q \not \equiv 0$ and $Q,P$ have no common factor, such that (up to rotation on $\Ss^2$) the rational function $\uu=(u_1, u_2, u_3)$ is given by the inverse stereographic projection applied to $R(x)$. Now, let us take sequences $P_n, Q_n \in \C[X]$ as provided in Lemma \ref{lem:rat_approx}. Define $R_n(x) = P_n(x)/Q_n(x) \in \C(X)$ and consider the sequence of rational maps $\uu^{(n)} : \R \to \Ss^2$ with components
$$
u_1^{(n)}(x) + \ii u_2^{(n)}(x)= \frac{2 R_n(x)}{|R_n(x)|^2 +1} , \quad u_3^{(n)}(x) = \frac{|R_n(x)|^2-1}{|R_n(x)|^2 +1} \, .
$$
Since $|P_n|^2 + |Q_n|^2$ has only simple zeros, we see that each rational map $\uu^{(n)}$ has only simple poles. Applying the known implication $(iii) \Rightarrow (ii)$, we conclude that $u_1^{(n)}, u_2^{(n)}, u_3^{(n)}$ for all $n$ have a least common denominator of degree $2N$. Thus for every rational map $\uu^{(n)} : \R \to \Ss^2$ we can apply Lemma \ref{lem:simple_poles} to conclude that
$$
\mathrm{Rank} \, (K_{\Ub_n}) = N \quad \mbox{for all $n \in \N$} \, .
$$
On the other hand, since $\uu^{(n)}(x) \to \uu(x)$ pointwise and $|\uu^{(n)}(x)| = 1$, we see that $K_{\Ub_n} f \to K_{\Ub} f$ in $L^2(\R, \C^2)$ for every $f \in L^2_+(\R, \C^2)$ by dominated convergence. From this we easily deduce that
$$
N=\liminf_{n \to \infty} \mathrm{Rank} (K_{\Ub_n}) \geq \mathrm{Rank}(K_{\Ub}) \, .
$$ 
This completes the proof that \eqref{eq:rank_K} holds whenever $u_1, u_2, u_3$ have a least common denominator of degree $2N$.

\medskip
\textbf{$(i) \Rightarrow (ii)$}.  Suppose now that $\mathrm{Rank}(K_{\Ub}) = N$ holds for some integer $N \geq 1$. Let $D \in \R[X]$ denote the least common denominator of $u_1, u_2, u_3$. Since $D$ has no zeros in $\R$, we must have that $\deg D = 2 m$ for some integer $m \geq 1$. We claim that 
$$
m=N \, .
$$ 
Indeed, if $\uu : \R \to \Ss^2$ has simple poles (in the sense of Lemma \ref{lem:simple_poles}), we can use Lemma \ref{lem:simple_poles} directly to deduce that $m=N$ must hold. 

If $\uu : \R \to \Ss^2$ has multiple poles, then $\deg D = 2m$ where $m \geq 1$ is the number poles $\uu$ counted with multiplicity. By the same argument using approximation with simple pole rational functions $\uu^{(n)} : \R \to \Ss^2$ as in the previous step, we conclude that $\mathrm{Rank}(K_{\Ub}) = m$. Hence $m=N$ is also true in this case.

The proof of Theorem \ref{thm:toeplitz_stereo_app} is now complete.
\end{proof}

\end{document}
